% Part: intuitionistic-logic
% Chapter: tableaux
% Section: proofs

\documentclass[../../../include/open-logic-section]{subfiles}

\begin{document}

\olfileid{int}{tab}{prf}

\olsection{د شهودي منطق لپاره \usetoken{P}{tableau}}

\begin{ex}
  داسې تړلے تابلو وړاندې کوو چې $(!A \land !B) \lif !C \Proves
  !A \lif (!B \lif !C)$ ښيي۔ په اوومه کرښه کښې نارښتينے
  عطف دوه څانګې جوړوي۔ په لومړۍ کښې له پخواني مخکښ څخه
  د عطف د لومړي غړي صدق پاتې کېږي، او په دويمه کښې په هماغه
  مخکښ دويم غړے رښتينے دے۔ په دريمه څانګه کښې پايله په يوه
  مخکښ هم رښتينې او هم نارښتينې ده۔
  \begin{oltableau}
    [\pFmla{\True}{(\formula{A} \land \formula{B}) \lif \formula{C}}{1},
      just =\TAss
      [\pFmla{\False}{\formula{A} \lif (\formula{B} \lif \formula{C})}{1},
        just =\TAss
        [\pFmla{\True}{\formula{A}}{1.1},
          just = {\TRule{\False}{\lif}[2]}
          [\pFmla{\False}{\formula{B} \lif \formula{C}}{1.1},
            just = {\TRule{\False}{\lif}[2]}
            [\pFmla{\True}{\formula{B}}{1.1.1},
              just = {\TRule{\False}{\lif}[4]}
              [\pFmla{\False}{\formula{C}}{1.1.1},
                just = {\TRule{\False}{\lif}[4]}
                [\pFmla{\False}{\formula{A} \land \formula{B}}{1.1.1},
                  just = {\TRule{\True}{\lif}[1]}
                  [\pFmla{\False}{\formula{A}}{1.1.1},
                    just= {\TRule{\False}{\land}[7]}, close]
                  [\pFmla{\False}{\formula{B}}{1.1.1},
                    just= {\TRule{\False}{\land}[7]}, close]]
                [\pFmla{\True}{\formula{C}}{1.1.1},
                  just = {\TRule{\True}{\lif}[1]}, close]]
            ]
          ]
        ]
      ]
    ]
  \end{oltableau}
\end{ex}



\begin{prob}
  د لاندې بيانونو د ښودلو لپاره تړلي شهودي !!{tableau}s ومومئ:
  \begin{enumerate}
    \item $\Proves !A \lif (!B \lif !A)$
    \item $\Proves \lnot(!A \land \lnot !A)$
    \item $!A \lif (!B \lif !C) \Proves (!A \land !B) \lif !C$
    \item $\lnot !A \lor \lnot !B \Proves \lnot(!A \land !B)$
  \end{enumerate}
\end{prob}

\end{document}
